\section{The finite proof kernel}
\label{sec:kernel}

This section states the reusable lemmas from which the compiler is built.
They are intentionally small.  Each either follows directly from the checked
natural-deduction calculus or reduces to quantifier-free induction and checked
arithmetic.

\subsection{Certificate sealing}

\begin{theorem}[Local certificate sealing]
\label{thm:sealing}
Let \(A(c,\vec x)\) and \(B(\vec x)\) be formulas, with \(c\) absent from
\(B\).  If
\[
 \Gamma\vdash_{\PRAAx}\exists c\,A(c,\vec x)
 \qquad\text{and}\qquad
 A(c,\vec x),\Gamma^+\vdash_{\PRAAx}B(\vec x)^+,
\]
then \(\Gamma\vdash_{\PRAAx}B(\vec x)\).
\end{theorem}

\begin{proof}
This is exactly the existential-elimination constructor \texttt{Prf.exE} of
\texttt{ConciseIPRA.Deduction}.  The shifts in the displayed second premise
express that the fresh certificate variable is in scope only in the branch.
No choice function is created.
\end{proof}

\begin{corollary}[Observational escape]
Suppose every valid certificate \(c\) gives a conclusion \(B(o)\) for
\(o=O(q,c)\), and valid certificates give equivalent observations.  If the
public conclusion is invariant under that equivalence, the certificate can be
sealed and only the invariant conclusion escapes.
\end{corollary}

\begin{proof}
Open the totality witness with \cref{thm:sealing}, derive the invariant
conclusion from its observation, and close the witness.  Equivalence is needed
only when two independently generated certificates are compared.
\end{proof}

\subsection{Simultaneous residual discharge}

\begin{theorem}[Finite residual discharge]
\label{thm:cut}
If \(\Gamma\vdash_{\PRAAx}\varphi\) and for every \(\psi\in\Gamma\) one has
\(\Delta\vdash_{\PRAAx}\psi\), then
\(\Delta\vdash_{\PRAAx}\varphi\).
\end{theorem}

\begin{proof}
This is the checked theorem \texttt{Prf.trans}.  Its proof repeatedly applies
implication introduction and modus ponens, and therefore requires no cut rule
as a new primitive.
\end{proof}

This theorem is the formal meaning of ``plugging producers into a compiled
proof''.  A framework report that lists residual assumptions but does not
produce this sequent has not yet specified composition at the object level.

\subsection{Booleanization of fixed bounded formulas}

For each fixed bounded formula \(A(\vec x)\), FIBER generates a primitive
recursive bit \([A](\vec x)\).  Atomic equalities and inequalities use the
checked equality/order bits.  Connectives use Boolean operations.  A bounded
existential is a finite disjunction fold, and a bounded universal is a finite
conjunction fold.

\begin{theorem}[Fixed-formula Boolean reflection schema]
\label{thm:boolean-reflection}
For every fixed bounded formula \(A(\vec x)\), there is a primitive-recursive
bit \([A]\) and a generated object derivation
\[
  []\vdash_{\PRAAx}\forall\vec x\,
       \bigl([A](\vec x)=1\leftrightarrow A(\vec x)\bigr).
\]
The construction is effective in the external syntax of \(A\); it does not
require an internal universal formula evaluator.
\end{theorem}

\begin{proof}[Human proof and induction accounting]
Proceed by structural recursion on the fixed syntax.

For equality and order atoms, use the existing checked reflection lemmas for
equality tests, monus order, and comparison bits.  Boolean conjunction,
disjunction, and implication are proved by the finite truth tables, each a
fixed collection of arithmetic cases.

For a bounded existential \(\exists i<N\,A(i,\vec x)\), define
\[
 e(0,\vec x)=0,\qquad
 e(k+1,\vec x)=\max(e(k,\vec x),[A](k,\vec x)).
\]
The formula used in induction is not the quantified formula.  It is a
quantifier-free state equation about \(e\), a candidate index, and the
reflected bit.  The forward direction from \(e(N)=1\) obtains the first stage
at which the bit becomes one by checked bounded search, then applies the
inductive reflection theorem for \(A\).  The reverse direction takes an
explicit existential witness \(i<N\), proves by a quantifier-free loop lemma
that its bit contributes to \(e(N)\), and concludes \(e(N)=1\).

For a bounded universal, define the conjunction fold
\[
 a(0,\vec x)=1,\qquad
 a(k+1,\vec x)=\min(a(k,\vec x),[A](k,\vec x)).
\]
From \(a(N)=1\) and an arbitrary \(i<N\), a quantifier-free indexing lemma
proves \([A](i,\vec x)=1\), then reflection yields \(A(i,\vec x)\).  Conversely,
under the universal hypothesis, quantifier-free induction proves the state
formula \(a(k)=1\) for \(k\le N\); the successor step instantiates the
universal hypothesis at \(k\).  The induction predicate itself remains the
atomic equation \(a(k)=1\), optionally guarded by the quantifier-free
inequality \(k\le N\).

All recursion symbols and formulas are fixed before the derivation is
constructed.  Thus no truth predicate and no universal evaluator is used.
\end{proof}

\begin{remark}
The repository already contains most arithmetic pieces of this proof, but the
complete generic generated reflection library is one of the central Lean
milestones.  The report claims the mathematical schema, not an already
compiled implementation.
\end{remark}

\subsection{Bounded least selection}

Let \(v(q,c)\) be a primitive-recursive bit.  The checked bounded-search module
finds the least index below a bound at which a chosen bit convention succeeds.
A small adapter converts between ``zero means success'' and ``one means
success''.

\begin{theorem}[Canonical bounded selection]
\label{thm:least}
There is a primitive-recursive function \(\mu_v(q,B)\) such that \(\ipra\)
derives
\[
 \bigl(\exists c<B\;v(q,c)=1\bigr)
 \longrightarrow
 \left(
 \begin{array}{l}
 \mu_v(q,B)<B,\\[-2pt]
 v(q,\mu_v(q,B))=1,\\[-2pt]
 \forall j<\mu_v(q,B)\;v(q,j)=0
 \end{array}
 \right).
\]
\end{theorem}

\begin{proof}
Apply the checked least-zero search to the complement bit
\(1\mathbin{\dot-}v\).  Its specification supplies an accepted index no later
than any supplied accepted witness and rules out all earlier successes.  The
remaining steps are Boolean arithmetic and the existing order lemmas.
\end{proof}

\begin{principle}[Use search only at an escape boundary]
If a certificate is used only in a proof branch, apply existential elimination.
Invoke \cref{thm:least} only when a downstream term or public theorem must name
a code.  This avoids unnecessary normalization burdens and prevents a local
existence proof from being mistaken for global choice.
\end{principle}

\subsection{Finite master certificates}

Suppose that for each \(i<N\) a bounded decidable verifier \(v(i,c)\) has a
witness below \(B(i)\).  A common informal step says ``choose all witnesses''.
In intuitionistic arithmetic this should be made explicit.

\begin{theorem}[Master-certificate theorem]
\label{thm:master}
Assume \(\ipra\) derives
\[
       \forall i<N\;\exists c<B(i)\;v(i,c)=1.
\]
Then there is a primitive-recursive sequence code \(C(N)\) for which \(\ipra\)
derives
\[
       \forall i<N\;
       \bigl(\seq(C(N),i)<B(i)\wedge
             v(i,\seq(C(N),i))=1\bigr).
\]
The construction uses no choice principle.
\end{theorem}

\begin{proof}
Define \(m(i)=\mu_{v(i,\cdot)}(B(i))\) by bounded least search.  The premise and
\cref{thm:least} prove the validity of \(m(i)\) for arbitrary \(i<N\).  Define
\(C(0)\) as the empty sequence and append \(m(k)\) at stage \(k+1\).  The
sequence equations are primitive recursive.  Quantifier-free induction on
\(k\) proves the decoding equations for entries already appended.  For an
arbitrary \(i<N\), instantiate the decoding lemma and the validity theorem for
\(m(i)\).  Mixed-radix coding gives an alternative when all component bounds
are available in advance.
\end{proof}

The theorem resolves hidden finite choice in proof graphs, finite covers,
strategy tables, and compatibility diagrams.  Notice that decidability and a
bound are essential; an unbounded family of arbitrary existential statements
would be stronger.

\subsection{Finite monotone demand closure}

Adaptive proofs can have cyclic precision requirements: node \(i\) asks for a
resource determined by the resources requested by other nodes.  Unbounded
fixed-point language is unnecessary when each coordinate has a cap.

\begin{theorem}[Least finite demand closure]
\label{thm:demand-closure}
Let
\[
      P=\prod_{i<r}\{0,1,\ldots,C_i\}
\]
with coordinatewise order, and let \(F:P\to P\) be monotone.  Define
\[
 d^0=0,\qquad d^{n+1}=d^n\vee F(d^n),
\]
where \(\vee\) is coordinatewise maximum.  Put \(K=\sum_{i<r}C_i\).  Then:

\begin{enumerate}
\item \(d^K=d^{K+1}\), so \(d^K\) is a prefixed point of \(F\);
\item if \(e\in P\) satisfies \(F(e)\le e\), then \(d^K\le e\);
\item consequently \(d^K\) is the least prefixed point of \(F\).
\end{enumerate}
\end{theorem}

\begin{proof}
The sequence is increasing.  If \(d^{n+1}\ne d^n\), at least one coordinate
increases by at least one, so the potential
\(\Phi(d)=\sum_{i<r}d_i\) strictly increases.  The potential starts at zero
and is at most \(K\).  Hence a strict increase cannot occur at all of the first
\(K+1\) stages, and monotonicity of the update makes the first equality
permanent.  This proves stabilization by stage \(K\).

For minimality, prove by induction on \(n\) that \(d^n\le e\).  It is true for
zero.  If \(d^n\le e\), monotonicity gives \(F(d^n)\le F(e)\le e\), and hence
\(d^{n+1}=d^n\vee F(d^n)\le e\).  In particular \(d^K\le e\).

Every operation is finite and primitive recursive.  The formal proof uses
encoded arrays, bounded sums, order, and quantifier-free induction on the
iteration counter.  The strict-increase argument can be expressed by the
potential equation and the checked order library.
\end{proof}

\subsection{Finite strategy certificates}

A fixed bounded quantifier alternation can be represented without classical
choice.

\begin{definition}[Strategy certificate]
For a fixed formula
\[
 Q_0x_0<B_0\;Q_1x_1<B_1(\vec x_{<1})\cdots
 Q_sx_s<B_s(\vec x_{<s})\;A(\vec x),
\]
where \(A\) is reflected and each \(Q_i\) is \(\forall\) or \(\exists\), a
strategy certificate contains, at every existential node, a response as a
function of the preceding universal moves, and enough finite table data to
verify every leaf.
\end{definition}

\begin{theorem}[Bounded strategy realization]
\label{thm:strategy}
For each fixed bounded quantifier pattern there is a primitive-recursive
strategy verifier and a generated \(\ipra\) derivation from ``the strategy
verifies'' to the bounded formula.
\end{theorem}

\begin{proof}
Induct externally on the fixed quantifier pattern.  At an existential node,
decode the response selected by the strategy, verify its bound, and apply
existential introduction.  At a universal node, introduce an arbitrary bounded
move and select the corresponding certified subtree.  At a leaf, apply
\cref{thm:boolean-reflection}.  The only induction internal to \(\ipra\) checks
finite table lookup and bounds, so its invariant is quantifier-free.  No law of
excluded middle is used beyond reflected decidability of the bounded atoms.
\end{proof}

The theorem covers finite games, bounded minimax arguments, finite adversarial
verification, and many bounded classical case analyses.  It does not turn an
unbounded classical proof into a witness-producing intuitionistic proof.
